\section{总结与延伸}

NLLB 的核心贡献可以压缩成一句话：\textbf{大规模多语翻译是一项人本需求、评测基础设施、数据生产、模型容量、训练调度与安全验收共同设计的系统工程。} 如果只记住 MoE 或 44\% BLEU，就会错过课程真正希望传达的方法论。

\subsection{十二条可执行结论}

\begingroup
\footnotesize
\setlength{\columnsep}{1.2em}
\begin{multicols}{2}
\begin{enumerate}[itemsep=0.08em,topsep=0.15em,parsep=0pt,leftmargin=*]
  \item 先确认社区需求，再决定语言、脚本、方向与场景。
  \item 把 high/low-resource 定义为数据条件，不按人口或价值排序。
  \item 在训练前建设可比、人工翻译、跨方向的 evaluation set。
  \item 为语言标准、译者、reviewer 和 post-editing 留出真实成本。
  \item 用少量可信 seed 校准 LID、encoder、filter 与初始 MT。
  \item 把 web mining 设计为可回溯的迭代 pipeline，而不是一次抓取。
  \item 对每份数据记录 provenance、noise、scale 与 teacher dependency。
  \item 长尾平衡要增加有效多样性，不能只复制低资源样本。
  \item MoE 的容量收益必须和低资源过拟合、路由与正则一起评估。
  \item 按 direction、resource group 和 metric 报告结果，不只报平均值。
  \item Automatic、human 与 toxicity evaluation 分别回答不同问题。
  \item 开放 benchmark、数据工具和模型，才能让社区继续修正覆盖缺口。
\end{enumerate}
\end{multicols}
\endgroup

\subsection{最终系统图景}

最后把整堂课串起来，可以得到一条带反馈的工程闭环。这个闭环不是线性瀑布流程，因为人类评估、安全失败与社区反馈都会迫使团队返回标准、数据和模型阶段：
\[
\begin{aligned}
\text{Community needs}
&\rightarrow\text{Standards \& Evaluation}
\rightarrow\text{Seed Data}
\rightarrow\text{Mining / BT}\\
&\rightarrow\text{MoE Training}
\rightarrow\text{Human \& Safety Validation}
\rightarrow\text{Community feedback}.
\end{aligned}
\]
每个箭头都可能引入选择偏差，每个反馈都可能改变前序标准。真正成熟的系统不是一次训练完成，而是能够解释数据从哪里来、模型为何在某些语言失败、哪些指标支持结论，以及谁有权修正错误。

\begin{importantbox}{最终判断标准}
“No Language Left Behind” 不是达到某个语言数后即可宣布完成。更可靠的判断是：尾部语言是否拥有可审计数据与评测；母语者是否参与需求、标准和验收；系统是否报告方向级质量与高风险错误；开放资产是否允许社区复现、批评和改进。覆盖数字只能作为这套证据中的一项。
\end{importantbox}

\subsection{拓展阅读}

\begingroup
\scriptsize
\begin{itemize}[itemsep=0pt,topsep=0.15em,parsep=0pt]
  \item NLLB Team, \href{https://arxiv.org/abs/2207.04672}{No Language Left Behind: Scaling Human-Centered Machine Translation}：完整的数据、模型、评测与社会影响报告。
  \item Meta AI, \href{https://research.facebook.com/publications/no-language-left-behind/}{NLLB publication page} 与 \href{https://github.com/facebookresearch/fairseq/tree/nllb}{fairseq NLLB release}：官方项目与模型代码入口。
  \item \href{https://github.com/facebookresearch/flores}{FLORES}：多语评测数据；\href{https://github.com/facebookresearch/stopes}{Stopes}：可扩展多语数据挖掘工具。
  \item Schwenk et al., \href{https://arxiv.org/abs/1907.05791}{WikiMatrix} 与 \href{https://github.com/facebookresearch/LASER}{LASER}：句向量语料挖掘基础。
\end{itemize}
\endgroup

\end{document}
